% !TEX encoding = UTF-8 Unicode
\documentclass[11pt,a4paper]{article}

\usepackage[utf8]{inputenc}
\usepackage[T1]{fontenc}
\usepackage{geometry}
\usepackage{hyperref}
\usepackage{enumitem}
\usepackage{booktabs}
\usepackage{array}
\usepackage{tabularx}
\usepackage{xcolor}
\usepackage{fancyhdr}
\usepackage{titlesec}
\usepackage{parskip}
\usepackage{listings}
\usepackage{tcolorbox}
\tcbuselibrary{breakable, skins, listings}

\geometry{a4paper, top=2.2cm, bottom=2.2cm, left=2.5cm, right=2.5cm}

\definecolor{accentblue}{RGB}{30, 80, 160}
\definecolor{promptbg}{RGB}{242, 244, 248}
\definecolor{promptborder}{RGB}{30, 80, 160}
\definecolor{warnbg}{RGB}{255, 248, 235}
\definecolor{warnborder}{RGB}{190, 80, 0}
\definecolor{tipbg}{RGB}{240, 248, 240}
\definecolor{tipborder}{RGB}{50, 130, 60}
\definecolor{darkgray}{RGB}{70, 70, 70}
\definecolor{medgray}{RGB}{130, 130, 130}
\definecolor{cheatbg}{RGB}{250, 250, 252}
\definecolor{cheatborder}{RGB}{30, 80, 160}
\definecolor{labelcolor}{RGB}{40, 40, 100}

\pagestyle{fancy}
\fancyhf{}
\fancyhead[L]{\small\textcolor{medgray}{Airbnb európai szálláshelyek elemzése}}
\fancyhead[R]{\small\textcolor{medgray}{Sapientia EMTE -- dr.\ Pál László}}
\fancyfoot[C]{\small\textcolor{medgray}{\thepage}}
\renewcommand{\headrulewidth}{0.3pt}

\titleformat{\section}{\large\bfseries\color{accentblue}}{\thesection}{0.6em}{}%
  [\vspace{-2pt}\textcolor{accentblue}{\rule{\linewidth}{0.6pt}}]
\setcounter{secnumdepth}{0}
\titleformat{\subsection}{\normalsize\bfseries\color{darkgray}}{}{0em}{}
\titlespacing{\section}{0pt}{18pt}{8pt}
\titlespacing{\subsection}{0pt}{12pt}{4pt}

\lstdefinestyle{promptlst}{
  basicstyle=\small\ttfamily\color{black},
  breaklines=true,
  breakatwhitespace=false,
  frame=none,
  columns=fullflexible,
  keepspaces=true,
  showstringspaces=false,
  literate=
    {á}{{\'{a}}}1 {é}{{\'{e}}}1 {í}{{\'{i}}}1 {ó}{{\'{o}}}1 {ö}{{\"o}}1
    {ő}{{\H{o}}}1 {ú}{{\'{u}}}1 {ü}{{\"u}}1 {ű}{{\H{u}}}1
    {Á}{{\'{A}}}1 {É}{{\'{E}}}1 {Í}{{\'{I}}}1 {Ó}{{\'{O}}}1 {Ö}{{\"O}}1
    {Ő}{{\H{O}}}1 {Ú}{{\'{U}}}1 {Ü}{{\"U}}1 {Ű}{{\H{U}}}1,
}

\tcbset{
  promptstyle/.style={
    breakable, enhanced,
    colback=promptbg, colframe=promptborder,
    boxrule=1pt, arc=3pt,
    left=10pt, right=10pt, top=10pt, bottom=8pt,
    attach boxed title to top left={xshift=10pt, yshift=-\tcboxedtitleheight/2},
    boxed title style={
      colback=promptborder, colframe=promptborder,
      arc=2pt, boxrule=0pt,
      left=6pt, right=6pt, top=2pt, bottom=2pt
    },
    coltitle=white, fonttitle=\footnotesize\bfseries\sffamily,
  },
  warnstyle/.style={
    breakable, colback=warnbg, colframe=warnborder,
    boxrule=0.5pt, leftrule=3pt, arc=2pt,
    left=8pt, right=8pt, top=5pt, bottom=5pt,
    fontupper=\small\sffamily,
  },
  tipstyle/.style={
    breakable, colback=tipbg, colframe=tipborder,
    boxrule=0.5pt, leftrule=3pt, arc=2pt,
    left=8pt, right=8pt, top=5pt, bottom=5pt,
    fontupper=\small\sffamily,
  },
  cheatstyle/.style={
    breakable, enhanced, colback=cheatbg, colframe=cheatborder,
    boxrule=0.8pt, arc=3pt,
    left=10pt, right=10pt, top=9pt, bottom=7pt,
    attach boxed title to top left={xshift=10pt, yshift=-\tcboxedtitleheight/2},
    boxed title style={
      colback=cheatborder, colframe=cheatborder,
      arc=2pt, boxrule=0pt,
      left=6pt, right=6pt, top=2pt, bottom=2pt
    },
    coltitle=white, fonttitle=\footnotesize\bfseries\sffamily,
  },
}

\newtcblisting{prompt}[1][Prompt]{
  promptstyle, title=#1,
  listing only,
  listing options={style=promptlst},
}

\newcommand{\warn}[1]{%
  \vspace{3pt}%
  \begin{tcolorbox}[warnstyle]%
  \textbf{\textcolor{warnborder}{Figyelem:}} #1%
  \end{tcolorbox}\vspace{2pt}}

\newcommand{\tip}[1]{%
  \vspace{3pt}%
  \begin{tcolorbox}[tipstyle]%
  \textbf{\textcolor{tipborder}{Tipp:}} #1%
  \end{tcolorbox}\vspace{2pt}}

\newcommand{\lbl}[1]{\vspace{6pt}\noindent{\small\bfseries\textcolor{labelcolor}{#1}}\par\smallskip}

\hypersetup{colorlinks=true, linkcolor=accentblue, urlcolor=accentblue}

\begin{document}

\thispagestyle{empty}
\begin{center}
  {\LARGE\bfseries Airbnb európai szálláshelyek elemzése}\\[5pt]
  {\normalsize Generatív AI és alkalmazások modul}\\[1pt]
  {\small\color{medgray} dr.\ Pál László, Sapientia EMTE, 2026. március 6--7.}
\end{center}

\vspace{10pt}
\noindent
\begin{tabularx}{\linewidth}{lXr}
\toprule
\textbf{Blokk} & \textbf{Tartalom} & \textbf{Időtartam} \\
\midrule
1.\ feladat    & Adatbetöltés & $\sim$10 perc \\
1.5\ feladat   & EDA (Exploratory Data Analysis) -- adatfeltáró elemzés & $\sim$10 perc \\
2--3.\ feladat & Alapstatisztikák, outlier szűrés, korreláció-heatmap & $\sim$20 perc \\
4--5.\ feladat & KDE, összefüggés-vizsgálatok & $\sim$25 perc \\
6.\ feladat    & K-means klaszterezés -- Budapest & $\sim$30 perc \\
7.\ feladat    & Ár-előrejelzés -- Random Forest & $\sim$20 perc \\
8.\ feladat    & Hasonló szállások ajánlója (opcionális) & $\sim$15 perc \\
\midrule
\multicolumn{2}{l}{\textbf{Teljes blokk}} & \textbf{$\sim$115 perc} \\
\bottomrule
\end{tabularx}

\smallskip
\noindent\textbf{Eszközök:} Google Colab + Gemini (\texttt{gemini.google.com})\\
\noindent\textbf{Dataset:} \texttt{airbnb\_europe.csv}

\vspace{8pt}

\begin{tcolorbox}[warnstyle]
\textbf{\textcolor{warnborder}{Oszlopnevek:}}
Amikor Gemini-nek promptot írsz, mindig pontosan ezeket az oszlopneveket add meg
(nagy/kis betű és szóközök számítanak):\\[4pt]
\texttt{Price}, \texttt{City}, \texttt{Day}, \texttt{Room Type},
\texttt{Superhost}, \texttt{Person Capacity}, \texttt{Cleanliness Rating},
\texttt{Guest Satisfaction}, \texttt{Bedrooms},
\texttt{City Center (km)}, \texttt{Metro Distance (km)},
\texttt{Normalised Attraction Index}
\end{tcolorbox}

%=============================================================
\section{1. feladat -- Adatbetöltés}
%=============================================================

\noindent\textit{Töltsd be az adatfájlt és vizsgáld meg az alapinformációkat.
Az \texttt{airbnb\_europe.csv} fájlt töltsd fel a Colab session-be
(bal oldali fájlkezelő -- feltöltés).}

\begin{prompt}[Prompt -- adatbetöltés]
Írj Python kódot Google Colab notebookban, amely:
- Betölti az airbnb_europe.csv fájlt pandas DataFrame-be (df)
- Kiírja a sorok és oszlopok számát
- Megjeleníti az első 5 sort (df.head())
- Kiírja az összes egyedi város nevét
- Ellenőrzi és kiírja a hiányzó értékeket oszloponként
\end{prompt}

%=============================================================
\section{1.5 feladat -- EDA (Exploratory Data Analysis) – adatfeltáró elemzés}
%=============================================================

\noindent\textit{Mielőtt bármit szűrünk vagy transzformálunk, nézzük meg, mit tartalmaz az adatbázis.
Ez a két diagram megmutatja, miért lesz szükség outlier szűrésre és log-transzformációra a következő lépésekben.}

\begin{prompt}[Prompt -- EDA vizualizáció]
Írj Python kódot, amely:
- Barplot: átlagos 'Price' városonként (sns.barplot, palette='Set3')
  Rendezd csökkenő sorrendbe az átlagár szerint
  Méret: 10x5, cím: 'Átlagos éjszakai ár városonként (EUR)'
  x tengelyen forgatás: 30 fok
- Hisztogram: a 'Price' oszlop nyers eloszlása (bins=60, color='steelblue')
  Méret: 10x4, cím: 'Ár eloszlása (nyers adatok)'
  x tengely felirat: 'Ár (EUR/éj)'

Meglévő változó: df (pandas DataFrame)
\end{prompt}

\tip{A barplotból látszik, melyik városok drágábbak (pl.\ Amsterdam, Párizs vs.\ Budapest).
A hisztogramból látszik az erős jobbra ferdülés -- néhány szálláshely extrém magas ára torzítja az egész képet.
Ez magyarázza, miért kell majd outlier szűrés és log-transzformáció.}

%=============================================================
\section{2. feladat -- Alapstatisztikák és outlier szűrés}
%=============================================================

\noindent\textit{Vizsgáljuk meg az ár eloszlását és a leíró statisztikákat.
Ha a \texttt{describe()} eredménye alapján extrém kiugró értékek láthatók,
szűrjük ki őket egy megfelelő percentilis határral.}

\begin{prompt}[Prompt -- outlier szűrés]
Írj Python kódot, amely:
- Kiírja a df.describe() eredményét (alapstatisztikák)
- Rajzolja meg a 'Price' oszlop hisztogramját
- Ha a hisztogram alapján szükséges, szűrje ki a kiugró
  értékeket a 'Price' oszlopból egy megfelelő percentilis határral
- Írja ki, hány sor maradt és mi volt a szűrési határ (EUR/éj)

Meglévő változó: df (pandas DataFrame)
Oszlopnév: 'Price'
\end{prompt}

\warn{Az outlier szűrés mértéke befolyásolja az összes
eredményt -- a klaszterközéppontokat, a modell pontosságát és az
ajánló kimenetét is.}

%=============================================================
\section{3. feladat -- Korreláció-heatmap}
%=============================================================

\noindent\textit{A heatmap megmutatja, melyik változók függnek össze egymással.
Az értékek --1 és +1 között mozognak: a kék cellák pozitív, a piros cellák
negatív összefüggést jeleznek.}

\begin{prompt}[Prompt -- korreláció-heatmap]
Írj Python kódot, amely elkészíti a korreláció-heatmapet:
- Csak numerikus oszlopokat szerepeltess (df.select_dtypes)
- sns.heatmap, annot=True, cmap='RdBu_r', vmin=-1, vmax=1
- linewidths=0.5, linecolor='black'
- Méret: 14x9
- Cím: 'Korrelációs mátrix'

Meglévő változó: df (szűrt DataFrame)
\end{prompt}

\tip{Figyeld meg a legerősebb összefüggéseket: \texttt{Cleanliness Rating}--\texttt{Guest Satisfaction} (0.69),
\texttt{Person Capacity}--\texttt{Bedrooms} (0.56).
Meglepetés: a \texttt{Price} és a \texttt{Guest Satisfaction} között nincs erős korreláció (0.02) --
a drágább szállás nem jelent jobb élményt.
Az \texttt{Attraction Index} és \texttt{Restaurant Index} változók egymással erősen korrelálnak (0.74--0.84):
ugyanazt a jelenséget mérik különböző normalizálással.
Ez a Random Forest feature importance-nál is látszani fog -- a ``szavazatok'' megoszlanak közöttük.}

%=============================================================
\section{4. feladat -- Ár eloszlás városonként (log-transzformáció)}
%=============================================================

\noindent\textit{A KDE (Kernel Density Estimation) egy simított hisztogram:
folytonos görbével mutatja meg, hol sűrűsödnek az árak.
Több város egyszerre jól összehasonlítható rajta.
A nyers ár eloszlása miatt azonban a városok görbéi alig különböztethetők meg --
néhány extrém drága szálláshely szétnyújtja a skálát.
A log10-transzformáció összenyomja a nagy értékeket,
így az eloszlások összehasonlíthatóvá válnak.
Például: 100 EUR $\to$ 2.0, 1000 EUR $\to$ 3.0.}

\begin{prompt}[Prompt -- log-transzformált KDE]
Írj Python kódot, amely:
- Létrehoz egy df_log másolatot a df-ből
- A df_log['Price'] oszlopot felváltja np.log10(df['Price'])-vel
- sns.displot KDE diagramot készít:
  x='Price', hue='City', kind='kde', fill=True
  height=6, aspect=1.8
- Cím: 'Ár eloszlás városonként (log10 skála)'
- x tengely felirat: 'log10(Ár)'

Meglévő változók: df, df_log
\end{prompt}

\tip{A görbék csúcsa mutatja a leggyakoribb árkategóriát városonként:
log10(2.0) = 100 EUR, log10(2.3) $\approx$ 200 EUR, log10(2.5) $\approx$ 316 EUR.
Rome csúcsa a legszűkebb és legmagasabb -- az árak kevésbé szóródnak.
Amsterdam és Paris görbéje jobbra tolódik és lapos -- drágább és változatosabb árszintű piac.
Athens és Budapest balra helyezkedik -- jellemzően olcsóbb szállások.
A 3.0 feletti értékek (1000 EUR+) minden városban ritkaságok,
de Amsterdam esetén a legsűrűbbek.}

%=============================================================
\section{5. feladat -- Összefüggés-vizsgálatok (binning)}
%=============================================================

\noindent\textit{Három kérdést vizsgálunk boxplot diagramokkal.
A binning technika a folytonos változókat (pl.\ távolság km-ben)
kategóriákra osztja, így boxplottal vizualizálhatók.}

\lbl{1.\ diagram -- Tisztaság és vendégelégedettség}

\begin{prompt}[Prompt -- cleanliness vs. elégedettség]
Írj Python kódot seaborn boxplothoz:
- x='Cleanliness Rating', y='Guest Satisfaction'
- palette='Set3', width=0.7, whis=5
- Méret: 12x5
- Cím: 'Tisztaság hatása a vendégelégedettségre'

Meglévő változó: df
Oszlopnevek: 'Cleanliness Rating', 'Guest Satisfaction'
\end{prompt}

\lbl{2.\ diagram -- Városközponttól való távolság és ár}

\begin{prompt}[Prompt -- távolság vs. ár]
Írj Python kódot, amely:
- pd.cut segítségével a 'City Center (km)' oszlopot
  kategóriákra osztja: 0-3, 3-6, 6-9, ..., 24-27 km
- Menti a df['dist_bin'] oszlopba
- sns.boxplot: x='dist_bin', y=np.log10(df['Price'])
  palette='Set3', width=0.7
- Méret: 12x5, x tengelyen 30 fokos forgatás
- Cím: 'Városközpont távolsága vs. ár (log10)'

Meglévő változók: df, df_log
Oszlopnevek: 'City Center (km)', 'Price'
\end{prompt}

\lbl{3.\ diagram -- Vendégelégedettség és ár}

\begin{prompt}[Prompt -- elégedettség vs. ár]
Írj Python kódot, amely:
- pd.cut segítségével a 'Guest Satisfaction' oszlopot
  10-es lépésekkel osztja kategóriákra: 0-10, 10-20, ..., 90-100
- Menti a df['sat_bin'] oszlopba
- sns.boxplot: x='sat_bin', y=np.log10(df['Price'])
  palette='Set3', width=0.7
- Méret: 12x5
- Cím: 'Vendégelégedettség vs. ár (log10)'

Meglévő változók: df, df_log
\end{prompt}

\tip{1.\ diagram: egyértelmű összefüggés -- a tisztább szálláshely magasabb vendégelégedettséggel jár.
2.\ diagram: az ár a városközponttól távolodva enyhe csökkenést mutat, de a különbség kisebb a vártnál --
log10(2.3) $\approx$ 200 EUR vs.\ log10(2.2) $\approx$ 158 EUR.
A 21--24 km-es kategória kilóg: valószínűleg repülőtér-közeli szállások emelik az átlagot.
3.\ diagram: a \texttt{Guest Satisfaction} és az ár között nincs érdemi összefüggés --
a drágább szállás nem jelent jobb értékelést, és fordítva sem.}

%=============================================================
\section{6. feladat -- K-means klaszterezés (Budapest)}
%=============================================================

\noindent\textit{A K-means tanítás nélküli (unsupervised) algoritmus: előre megcímkézett adat nélkül fedez fel
csoportokat a hasonlóságok alapján. Egyszerre veszi figyelembe az összes változót -- ár, távolság,
elégedettség, hálószobák száma stb. -- így üzleti szegmensek azonosíthatók (pl.\ budget, prémium,
nagy kapacitású). A folytonos változókat standardizálni kell, mert eltérő skálán vannak;
a kategorikus oszlopokat one-hot encodinggal bináris (0/1) értékekké alakítjuk.}

\subsection{6.1 -- Adatelőkészítés és elbow method}

\begin{prompt}[Prompt -- Budapest szűrés és elbow]
Írj Python kódot, amely:
- Szűri a budapesti adatokat: df_bp = df[df['City'] == 'Budapest'].copy()
- One-hot encoding a kategorikus oszlopokra (pd.get_dummies)
- Bool oszlopokat int8-ra konvertál
- StandardScaler-rel standardizálja a numerikus oszlopokat (X_scaled)
- Elbow method: KMeans K=2..15, inertia értékeket gyüjt és kirajzolja
  x='k', y='Inertia', marker='o'
- Cím: 'Elbow method -- optimális K meghatározása'

Meglévő változó: df
\end{prompt}

\warn{Ennél a görbénél nincs éles törés -- ez tipikus. K=4 és K=6 között érdemes kísérletezni:
futtasd le mindkettőt és azt válaszd, amelyik klaszterprofil üzletileg értelmesebb
(pl.\ budget, prémium, nagy kapacitású szegmens).}

\subsection{6.2 -- Klaszterezés és PCA vizualizáció}

\noindent\textit{A PCA (Principal Component Analysis) az összes változó információját sűríti 2 dimenzióba,
hogy vizualizálható legyen. A tengelyértékek önmagukban nem értelmezhetők -- csak a pontok
egymáshoz viszonyított távolsága számít: a közel lévő pontok hasonló jellemzőjű szállásokat jelentenek.
Az a lényeg, hogy a \textbf{színek szétválnak-e}: ez jelzi, hogy az algoritmus valódi csoportokat talált.}

\begin{prompt}[Prompt -- K-means és PCA]
Írj Python kódot, amely:
- KMeans(n_clusters=5, n_init='auto', random_state=42) futtat X_scaled-en
- A klaszter labeleket hozzáadja df_bp['cluster'] oszlopként
- Kiírja a klaszter profilokat: csoportonkénti átlagok
  'Price', 'Guest Satisfaction', 'City Center (km)', 'Bedrooms' oszlopokra
- PCA(n_components=2) projekció, scatter plot klaszterszín szerint
  (cmap='tab10', s=20, alpha=0.7)
- Cím: 'Budapesti szállások klaszterei (PCA projekció)'

Meglévő változók: df_bp, X_scaled
\end{prompt}

\tip{A profilok üzleti szegmenseket mutatnak: a legdrágább, legjobb elhelyezkedésű klaszter a prémium szegmens,
a magas távolság + alacsony ár kombináció a budget külvárosi szegmens,
a nagy tömegű középső klaszterek különböző árszintű átlagos szállásokat jelentenek.
A PCA diagramon a színek szétválása megerősíti, hogy az algoritmus valódi csoportokat talált.}

%=============================================================
\section{7. feladat -- Ár-előrejelzés (Random Forest, Budapest)}
%=============================================================

\noindent\textit{A Random Forest döntési fák sokaságából álló ensemble modell.
Most arra tanítjuk, hogy a szálláshely jellemzői alapján megjósolja az árat.
A feature importance megmutatja, melyik változó befolyásolja leginkább az árat.}

\subsection{7.1 -- Modell tanítása és értékelése}

\begin{prompt}[Prompt -- Random Forest modell]
Írj Python kódot, amely:
- df_bp-ből a célváltozó: y = df_bp['Price']
- Feature-ök: X = df_bp összes többi numerikus oszlopa
- Train/test split: 80/20, random_state=42
- RandomForestRegressor(n_estimators=500, random_state=42, n_jobs=-1)
- Kiírja Train és Test MAE, RMSE, R2 értékeket
- Top 15 feature importance vízszintes barplot
  Cím: 'Top 15 -- mi határozza meg az árat Budapesten?'

Meglévő változók: df_bp (one-hot kódolt DataFrame)
\end{prompt}

\tip{Test R²=0.60: a modell az árvarianciának 60\%-át magyarázza --
a maradék 40\% olyan tényezőktől függ, amelyek nincsenek az adatban
(fotók minősége, leírás stílusa, gazda válaszkészsége).
A Train/Test R² különbség (0.91 vs.\ 0.60) túlilleszkedést jelez.

Feature importance: az elhelyezkedés (\texttt{Metro Distance}, \texttt{City Center}) dominálja az árat --
ez összhangban van az 5.\ feladat boxplot eredményeivel.
Az \texttt{Attraction} és \texttt{Restaurant} indexek együtt $\sim$40\%-ot tesznek ki:
pontosan a 3.\ feladatnál jelzett szavazatmegoszlás.
Az EDA-ban látott összefüggések tehát visszaigazolódnak a gépi tanulási modellben.}

%=============================================================
\section{8. feladat -- Hasonló szállások ajánlója \textnormal{\small(opcionális)}}
%=============================================================

\noindent\textit{Ez a feladat opcionális -- ha az idő engedi, vagy házi projektként. Elvégzése nem feltétele a többi feladatnak.}

\smallskip
\noindent A \textbf{cosine similarity} két vektor közötti szög koszinuszát méri.
Ha két szálláshely numerikus jellemzői hasonlók,
a köztük lévő szög kicsi -- a hasonlóság értéke 1-hez közelít.
Az ajánló az összes városra dolgozik, így előfordulhat,
hogy egy budapesti szálláshely mellé bécsi kerül -- ez nem hiba,
hanem az azonos ár-érték profil következménye.

\begin{tcolorbox}[tipstyle]
\textbf{\textcolor{tipborder}{Továbbfejlesztés -- beágyazásos ajánló (embedding-alapú):}}
A mostani megközelítés numerikus jellemzőkön dolgozik (ár, távolság stb.).
A modern ajánlórendszerek ennél tovább mennek: a szálláshely szöveges leírásából
nyelvi modell segítségével \textbf{beágyazási vektorokat} (embeddings) generálnak,
és a szemantikai hasonlóságot mérik -- nem csak a számokat.
Ehhez pl.\ a \texttt{sentence-transformers} könyvtár használható Colabon,
API kulcs nélkül (\texttt{paraphrase-multilingual-MiniLM} modell).
Ez a házi projekt egyik lehetséges iránya.
\end{tcolorbox}

\begin{prompt}[Prompt -- cosine similarity ajánló]
Írj Python kódot, amely felépíti a szállásajánlót:
- df-ből válaszd ki ezeket a feature-öket:
  'Price', 'Guest Satisfaction', 'City Center (km)',
  'Bedrooms', 'Normalised Attraction Index', 'Metro Distance (km)'
- StandardScaler standardizálás
- Írj ajanlj(idx, top_n=5) függvényt, amely:
  - cosine_similarity segítségével megkeresi a top_n
    leghasonlóbb szállást (saját magát kizárva)
  - Kiírja a referencia szállást és a top 5 ajánlást:
    'City', 'Room Type', 'Price', 'City Center (km)',
    'Guest Satisfaction', hasonlóság értéke
- Teszteld egy budapesti és egy párizsi szállásra

Szükséges import: from sklearn.metrics.pairwise import cosine_similarity
Meglévő változó: df (teljes, 10 városos DataFrame)
\end{prompt}

\tip{Az ajánló városok között is keres -- egy budapesti prémium szállás mellé bécsi kerülhet,
mert az ár-érték profiljuk hasonló. Ez a módszer korlátja is: csak a numerikus jellemzőket látja,
a város karakterét, a nyelvet vagy a kulturális kontextust nem.
Pontosan ezért mutat tovább az embedding-alapú megközelítés.}

%=============================================================
\newpage
\section*{Cheat Sheet -- Airbnb elemzés}
%=============================================================

\begin{tcolorbox}[cheatstyle, title={Adatbetöltés és tisztítás}]
\begin{tabularx}{\linewidth}{@{}lX@{}}
\texttt{pd.read\_csv('airbnb\_europe.csv')} & fájl betöltése \\[2pt]
\texttt{df.isnull().sum()} & hiányzó értékek \\[2pt]
\texttt{df['Price'].quantile(0.99)} & 99.\ percentilis \\[2pt]
\texttt{df[df['Price'] <= p99]} & outlier szűrés \\[2pt]
\texttt{np.log10(df['Price'])} & log-transzformáció \\[2pt]
\texttt{pd.cut(df[x], bins=..., labels=...)} & binning \\
\end{tabularx}
\end{tcolorbox}

\smallskip

\begin{tcolorbox}[cheatstyle, title={K-means klaszterezés -- lépések}]
\begin{tabularx}{\linewidth}{@{}lX@{}}
\textbf{1.} & Kategorikus oszlopok: \texttt{pd.get\_dummies} \\[2pt]
\textbf{2.} & Standardizálás: \texttt{StandardScaler().fit\_transform(X)} \\[2pt]
\textbf{3.} & Elbow: K=2..15, inertia görbe, keresd a töréspontot \\[2pt]
\textbf{4.} & \texttt{KMeans(n\_clusters=K, n\_init='auto', random\_state=42)} \\[2pt]
\textbf{5.} & Profilok: \texttt{df.groupby('cluster')[features].mean()} \\[2pt]
\textbf{6.} & PCA vizualizáció: \texttt{PCA(n\_components=2)} \\
\end{tabularx}
\end{tcolorbox}

\smallskip

\begin{tcolorbox}[cheatstyle, title={Random Forest -- lépések}]
\begin{tabularx}{\linewidth}{@{}lX@{}}
\textbf{1.} & \texttt{train\_test\_split(X, y, test\_size=0.2)} \\[2pt]
\textbf{2.} & \texttt{RandomForestRegressor(n\_estimators=500, n\_jobs=-1)} \\[2pt]
\textbf{3.} & Metrikák: MAE, RMSE, R2 -- train és test halmazon \\[2pt]
\textbf{4.} & \texttt{rf.feature\_importances\_} -- melyik változó számít? \\
\end{tabularx}
\end{tcolorbox}

\smallskip

\begin{tcolorbox}[cheatstyle, title={Cosine similarity ajánló (opcionális)}]
\begin{tabularx}{\linewidth}{@{}lX@{}}
\textbf{1.} & Feature-ök standardizálása \\[2pt]
\textbf{2.} & \texttt{cosine\_similarity(matrix[i].reshape(1,-1), matrix)} \\[2pt]
\textbf{3.} & Saját magát kizárjuk: \texttt{sim[idx] = 0} \\[2pt]
\textbf{4.} & Top N: \texttt{sim.argsort()[::-1][:n]} \\[2pt]
\textbf{+} & Továbbfejlesztés: \texttt{sentence-transformers} szöveges beágyazással \\
\end{tabularx}
\end{tcolorbox}

\smallskip

\begin{tcolorbox}[cheatstyle, title={Vibe coding -- hasznos promptminták}]
\begin{tabularx}{\linewidth}{@{}lX@{}}
\textit{Adatbetöltés} & \textit{``Írj kódot, amely betölti a CSV-t és kiírja az alapinformációkat''} \\[3pt]
\textit{Vizualizáció} & \textit{``Adj hozzá tengelyfeliratot, címet és rácsozást''} \\[3pt]
\textit{Hiba esetén} & \textit{``Ez a hibaüzenet jelent meg: [...] Javítsd ki''} \\[3pt]
\textit{Finomítás} & \textit{``A diagram túl zsúfolt, csökkentsd a betűméretet''} \\[3pt]
\textit{Magyarázat} & \textit{``Magyarázd el mit csinál ez a kódrészlet''} \\
\end{tabularx}
\end{tcolorbox}

\vfill
\begin{center}
\small\textcolor{medgray}{Sapientia EMTE \textbullet{} Generatív AI és üzleti alkalmazások \textbullet{} 2026}
\end{center}

\end{document}